% =============================================================================
% Article-track manuscript scaffold (first paper).
%
% Branch:        repo-publication-readiness
% Status:        SCAFFOLD ONLY -- experimental results not yet inserted.
% Methodology:   article-track (true N-objective NBI, VRF/FMSE, MBPA),
%                NOT identical to the dissertation baseline. See
%                ../docs/METHODOLOGY_DECISIONS.md and
%                DISSERTATION_TO_ARTICLE_MAP.md.
% Compile:       pdflatex main && bibtex main && pdflatex main && pdflatex main
%
% Notes
% - Uses a portable preamble (article class + natbib + hyperref) so the
%   scaffold compiles without journal-specific style files. Switch to a
%   target template (Elsevier elsarticle, IEEE, Springer Nature) once the
%   target venue is selected.
% - Every section file contains explicit TODO blocks for places where
%   results, citations, or experiment-dependent prose are still missing.
% =============================================================================

\documentclass[11pt,a4paper]{article}

% --- Packages -----------------------------------------------------------------
\usepackage[utf8]{inputenc}
\usepackage[T1]{fontenc}
\usepackage{lmodern}
\usepackage[english]{babel}
\usepackage{amsmath,amssymb,amsthm,bm}
\usepackage{graphicx}
\usepackage{booktabs}
\usepackage{multirow}
\usepackage{caption}
\usepackage{subcaption}
\usepackage{microtype}
\usepackage{xcolor}
\usepackage{enumitem}
\usepackage[round]{natbib}
\usepackage[colorlinks=true,linkcolor=blue!50!black,citecolor=blue!50!black,
            urlcolor=blue!50!black]{hyperref}
\usepackage{cleveref}
\usepackage{listings}

% --- Macros -------------------------------------------------------------------
\newcommand{\todoblock}[1]{\par\noindent\fbox{\parbox{0.97\linewidth}{%
  \textbf{TODO:} #1}}\par\smallskip}
\newcommand{\NBI}{NBI}
\newcommand{\VRF}{VRF}
\newcommand{\MBPA}{MBPA}
\newcommand{\CCDFC}{CCDFC}

\graphicspath{{figures/}}

% --- Title --------------------------------------------------------------------
\title{Multi-objective Hyperparameter Optimization of\\
       Gradient-Boosted Decision Trees via\\
       Design of Experiments and True Normal Boundary Intersection\\
       \small (working title)}

\author{%
  Caio Tertuliano Ribeiro\thanks{Federal University of Itajub\'{a} (UNIFEI).
    Corresponding author: \texttt{caio.tertu99@gmail.com}.}\\
  \textit{(co-authors TBD)}\\[0.4em]
}
\date{\today}

% =============================================================================
\begin{document}

\maketitle

\begin{abstract}
% Abstract: tight, claim-light. Numbers go in only after the v1
% experiment runs. Until then, keep methodological framing.
Hyperparameter tuning of gradient-boosted decision trees is inherently
multi-objective: predictive quality must be balanced against
computational cost and, in practice, against robustness across data
splits. We present a design-aware framework that integrates a Design of
Experiments (DoE) phase, design-aware response surface modeling
(process quadratic on factorial and central composite designs;
Scheff\'{e} canonical mixture polynomials on simplex weight designs), a
flexible factor-analysis layer (PCA with Varimax rotation, configurable
between automatic, fixed, manual-construct, and disabled modes), and a
true \emph{$N$-objective} Normal Boundary Intersection (NBI) routine
implementing the Das--Dennis sub-problem with explicit anchors, payoff
matrix, quasi-normal direction, and equality constraints rather than a
weighted-sum scalarization. A conditional Mixture-Based Performance
Assessment (MBPA) post-optimization stage refines weights when the
primary frontier is flat or weakly informative. We instantiate the
framework on the
$12{\times}3{\times}10$ benchmark introduced in
Section~\ref{sec:experimental-design} (twelve public tabular datasets,
three GBDT families, ten replicas under matched evaluation budgets) and
position it as a methodological refinement and generalization of the
dissertation pipeline of \citep{ribeiro2026dissertation}, anchored in
the multivariate NBI/MBPA formulation of \citep{pereira2025eaai}.
\todoblock{Insert quantitative claims (median CPU-hour reduction,
predictive parity bands, MBPA trigger rates) once the v1 experiment is
aggregated; write each claim verbatim from the result CSVs to avoid
drift.}

\end{abstract}

\noindent\textbf{Keywords:} multi-objective hyperparameter optimization,
design of experiments, response surface methodology, normal boundary
intersection, factor analysis, varimax rotation, mixture design,
gradient-boosted decision trees, XGBoost, LightGBM, CatBoost.

% --- Sections ----------------------------------------------------------------
\section{Introduction}\label{sec:introduction}

Hyperparameter optimization (HPO) is the dominant lever between an
out-of-the-box gradient-boosted decision tree (GBDT) implementation
and a competitive deployed model
\citep{probst2019tunability,bergstra2012random}. In tabular
classification, where GBDT methods such as XGBoost
\citep{chen2016xgboost}, LightGBM \citep{ke2017lightgbm}, and CatBoost
\citep{prokhorenkova2018catboost} remain the leading family
\citep{grinsztajn2022tree}, this lever simultaneously affects three
deliverables a practitioner cares about: predictive quality, the
wall-clock and energy cost of training, and the run-to-run stability
of the resulting model. Treating HPO as a single-objective scalar
search obscures these trade-offs and forces them to be re-discovered
post hoc.

A natural alternative is to formulate HPO as a multi-objective problem
and to seek a set of compromise solutions on the Pareto frontier
\citep{deb2002nsga2,morales2022survey}. Within this view, three
methodological challenges interact: (i) deciding \emph{which}
metrics enter the objective vector and \emph{in which direction} each
should be optimized; (ii) producing a sample of high-quality candidates
under a strict evaluation budget; and (iii) selecting a final
configuration in a way that is auditable rather than ad hoc.

The dissertation underlying this article \citep{ribeiro2026dissertation}
addressed (i)--(iii) with a Design of Experiments (DoE) front end, a
PCA + Varimax factor reduction of the response space (the
\emph{Varimax Rotated Factor Score} representation, $\VRF$), a
quadratic Response Surface Methodology (RSM) surrogate, and a
weighted-sum scalarization sweep over a $\beta$-grid that the
dissertation labelled ``NBI''. The post-defense audit summarized in
\texttt{docs/METHODOLOGY\_DECISIONS.md} of the companion repository
identified that the implementation was, in fact, a normalized
weighted-sum scalarization rather than the Das--Dennis Normal Boundary
Intersection \citep{das1998nbi}. Independently, the peer-reviewed
multivariate $\NBI$ formulation of \citet{pereira2025eaai} reframes
the same building blocks (rotated factor scores; payoff matrix;
mixture-based post-optimization on the simplex of weights) and
demonstrates the value of a true sub-problem-based $\NBI$ when
objectives are correlated.

This article presents the post-dissertation methodological refinement
of that pipeline as implemented on the
\texttt{repo-publication-readiness} branch of the companion repository.
The contributions are:
\begin{itemize}[leftmargin=1.2em]
  \item \textbf{Explicit objective directions and FMSE wrapping.}
        Every objective entering the optimization carries a required
        direction (\textsc{minimize}, \textsc{maximize},
        \textsc{target}) and an optional FMSE transform that converts
        target-seeking objectives to a uniform minimization convention,
        following Eqs.~21--22 of \citet{pereira2025eaai}.
  \item \textbf{True $N$-objective NBI.} Anchors are computed by
        per-objective surrogate optimization; the payoff matrix
        $\bm{\Phi}$ and the quasi-normal direction
        $\hat{\bm n} = -\bm{\Phi}\bm{1}/\|\bm{\Phi}\bm{1}\|$ are
        first-class objects; each sub-problem
        $\max\,t\;\text{s.t.}\;\hat{\bm F}(\bm x) = \bm F^{U} +
        \bm{\Phi}\bm\beta + t\,\hat{\bm n}$ is solved with vector
        equality constraints and SLSQP. The mathematical core is
        dimension-agnostic, supporting $q\ge 2$ objectives with no
        bi-objective hard-coding.
  \item \textbf{Design-aware surrogate modeling.} Process-quadratic
        RSM is used on factorial / central composite designs;
        Scheff\'{e} canonical mixture polynomials are used on
        simplex-lattice designs and the post-optimization stage. Cross
        usage is guarded by a validator that warns or rejects
        incompatible combinations.
  \item \textbf{Conditional MBPA post-optimization.} A second stage
        triggers only when the primary frontier diagnostics indicate a
        flat or compressed front; the trigger decision is logged
        regardless. When fired, it fits Scheff\'{e} mixture
        surrogates of generalized distance to utopia and Shannon
        entropy over the simplex of weights and refines $\bm w^{*}$
        inside an elliptical interior constraint.
  \item \textbf{Reproducibility and cost discipline.} A YAML config
        layer (Pydantic v2) makes every run declarative; a calibrated
        cost estimator predicts CPU-hours, wall-clock time, and cloud
        spend before committing to a campaign; a CI pipeline runs the
        $N$-objective NBI residual checks on every push.
\end{itemize}

The first article-track experiment, described in
Section~\ref{sec:experimental-design}, evaluates the framework on
twelve public tabular classification datasets, three GBDT algorithms,
and ten independent replicas, against four established baselines and a
legacy weighted-sum ablation. We deliberately defer the
82-dataset doctoral benchmark to a follow-up paper.

\paragraph{Outline.} Section~\ref{sec:related-work} surveys the
relevant HPO literature and positions this work against
\citet{pereira2025eaai} and \citet{ribeiro2026dissertation}.
Section~\ref{sec:method} develops the article-track method.
Section~\ref{sec:experimental-design} states the planned protocol.
Section~\ref{sec:results} reports results
\todoblock{after the v1 campaign runs}.
Section~\ref{sec:discussion} interprets them and
Section~\ref{sec:conclusion} closes.

\section{Related work}\label{sec:related-work}

We position the proposed method against six families of hyperparameter
optimization (HPO) algorithms, plus the OpenML benchmarking ecosystem
that defines our evaluation protocol. The selected baselines span
single-objective HPO under a fixed evaluation budget, multi-fidelity
optimization, evolutionary search, and multi-objective HPO; AutoML
systems are referenced as context rather than as direct baselines.

\subsection{Classical HPO baselines}\label{sec:rw-classical}
The dominant single-objective baselines for tabular HPO are random
search \citep{bergstra2012random}, Gaussian-process Bayesian
optimization (GP-BO) \citep{snoek2012bayesian,shahriari2016bayesian,
frazier2018tutorial}, and the tree-structured Parzen estimator (TPE)
\citep{bergstra2011tpe}, the latter as implemented by Hyperopt and
exposed in the form of the Optuna \texttt{TPESampler}
\citep{akiba2019optuna}. \citet{probst2019tunability} report that for
gradient-boosting decision trees (GBDT) most of the variance in
generalization is captured by a small set of tuned hyperparameters,
which makes random search a non-trivial control even at modest
budgets.

\subsection{Sequential model-based algorithm
configuration}\label{sec:rw-smac}
The Sequential Model-Based Algorithm Configuration (SMAC) family
\citep{hutter2011smac,lindauer2022smac3} replaces the GP surrogate
with a random forest, which natively handles mixed-type and
conditional hyperparameter spaces and tolerates noisy evaluations via
intensification and racing. SMAC3 is the de facto state of the art
for general-purpose HPO of tabular ML pipelines, and powers the
classical Auto-sklearn \citep{feurer2015autosklearn} backend.

\subsection{Multi-fidelity HPO}\label{sec:rw-multifidelity}
Multi-fidelity methods reuse cheap, low-fidelity evaluations
(e.g., few epochs, few CV folds, sub-sampled training set) to
prune unpromising configurations early. Successive Halving and
Hyperband \citep{li2017hyperband} formalize this as a bandit problem;
ASHA \citep{li2020asha} parallelizes the schedule for asynchronous
workers. BOHB \citep{falkner2018bohb} couples Hyperband's bracket
schedule with TPE-style model-based sampling, and DEHB
\citep{awad2021dehb} replaces TPE with differential evolution
\citep{storn1997differential} for better high-dimensional behavior.
For GBDT models, the most natural fidelity dimension is the number of
boosting iterations (\texttt{n\_estimators} for XGBoost / LightGBM /
CatBoost), which lets multi-fidelity methods reuse partially trained
ensembles cheaply.

\subsection{Evolutionary and derivative-free
HPO}\label{sec:rw-evolutionary}
Differential evolution \citep{storn1997differential} and
covariance-matrix adaptation evolution strategy (CMA-ES)
\citep{hansen2001cmaes} are derivative-free baselines for continuous
HPO; they are particularly relevant when the objective surface is
non-smooth or when the budget allows a generation-based search rather
than the strictly sequential single-incumbent dynamics of GP-BO. The
Nevergrad platform \citep{rapin2018nevergrad} is referenced as a
unified implementation harness for evolutionary and derivative-free
optimizers; the present article cites it for context rather than
reporting Nevergrad-only baselines.

\subsection{Multi-objective HPO}\label{sec:rw-moo}
\citet{morales2022survey} survey multi-objective HPO and group the
literature into evolutionary, scalarization-based, and decomposition
approaches. NSGA-II \citep{deb2002nsga2} and its many-objective
extension NSGA-III \citep{deb2014nsga3} are the canonical
evolutionary baselines (the latter is relevant for $q\!>\!2$
objectives via reference points, paralleling the simplex-lattice
weight grid we use for NBI). The multi-objective TPE variant MOTPE
\citep{ozaki2020motpe} is the natural multi-objective counterpart of
TPE-Hyperopt. ParEGO \citep{knowles2006parego} represents
scalarization-based BO with rotating weight vectors. Within the
DoE / RSM tradition, multi-objective problems are commonly attacked
with Normal Boundary Intersection \citep{das1998nbi}, which generates
evenly spaced points on the Pareto front of nonlinear MOO problems
and is well-defined for arbitrary $q$ objectives. The pymoo library
\citep{blank2020pymoo} provides reference Python implementations of
NSGA-II / NSGA-III and is the natural implementation candidate for the
evolutionary baselines in the comparative protocol.

\subsection{Multivariate NBI with factor
scores}\label{sec:rw-pereira}
The closest methodological reference for the article-track is
\citet{pereira2025eaai}, who replace the original responses with
Varimax-rotated factor scores ($\VRF$), wrap each $\VRF$ in an FMSE
target-seeking loss, and solve the resulting NBI sub-problem under a
spherical region constraint. They further introduce a
mixture-design-based post-optimization that fits Scheff\'{e}
canonical polynomials over the simplex of NBI weights and refines
$\bm w^{*}$ by minimizing a generalized distance to utopia subject to
an elliptical interior constraint and an entropy preference.

\subsection{AutoML systems as context}\label{sec:rw-automl}
Auto-sklearn \citep{feurer2015autosklearn}, FLAML
\citep{wang2021flaml} and AutoGluon-Tabular
\citep{erickson2020autogluon} bundle HPO with model selection,
ensembling, and resource control. They are cited as context: the
proposed method targets the HPO sub-problem of a fixed GBDT
algorithm, not full pipeline AutoML. FLAML is the most natural
practical baseline of the three (single-algorithm tuning with a CPU
budget) and is listed in the comparative protocol as an optional
addition; Auto-sklearn and AutoGluon are excluded from the headline
comparison because their search spaces are not directly aligned with
the per-algorithm hyperparameter list of XGBoost / LightGBM /
CatBoost.

\subsection{Benchmark suites}\label{sec:rw-benchmarks}
We use the OpenML-CC18 curated classification benchmark suite
\citep{bischl2021openmlbenchmark} as the primary evaluation panel.
OpenML-CC18 (\texttt{suite\_id = 99}) collects 72 standardized
classification tasks selected for license clarity, dataset diversity,
and well-defined train/test splits, and is widely used as a fair
common ground for HPO and AutoML evaluation. The OpenML platform
\citep{vanschoren2014openml} provides task identifiers, default folds,
and metadata that make task-based (rather than dataset-based)
benchmarks reproducible across implementations. Cross-method
statistical comparison follows \citet{demsar2006statistical}.

\subsection{This work in context}\label{sec:rw-positioning}
Relative to \citep{pereira2025eaai}, we (a) port the framework to HPO
of GBDT models with explicit hyperparameter box constraints rather
than a hyperspherical region; (b) keep the FMSE wrapping for $\VRF$
objectives but support a parallel direct-objective mode for raw ML
metrics with explicit directions; (c) make the NBI sub-problem the
\emph{actual} solver in the implementation rather than only a
mathematical reference. Relative to
\citep{ribeiro2026dissertation}, we replace the legacy weighted-sum
scalarization with the true $N$-objective NBI, generalize the
2-VRF / 3-VRF collapse to a configurable factor model
(\textsc{auto}/\textsc{fixed}/\textsc{manual}/\textsc{none}), make
RSM design-aware (separating process-quadratic from Scheff\'{e}
mixture models), and add the conditional MBPA stage. The legacy
weighted-sum solver is preserved as an explicit ablation baseline and
is never referred to as NBI in the article-track code or text. The
formal comparative protocol against the baselines listed above is
specified in \texttt{docs/COMPARATIVE\_PROTOCOL.md} and the
machine-readable method matrix at
\texttt{benchmarks/doctoral/openml\_cc18/method\_matrix.csv}.

\section{Method}\label{sec:method}

The article-track pipeline has six stages. Each stage corresponds to a
self-contained module in the companion repository (the module name is
indicated parenthetically below).

\subsection{Stage 1 -- Design of Experiments}
\label{sec:method-doe}
A user-supplied or library-generated experimental design
(\texttt{doe\_xgb.design}) covers the hyperparameter space. The
canonical artifact is the Minitab face-centered central composite
design ($\CCDFC$) of the dissertation, with $k{=}7$ XGBoost
hyperparameters, $n{=}88$ runs and four center points; the
\texttt{DesignProvider} also generates fractional factorial,
Box--Behnken, Latin hypercube, Sobol, and simplex-lattice designs
in pure NumPy. Each artifact carries matrix diagnostics (rank,
condition number, coverage of the coded box) and a SHA-256 of the
source CSV so reviewers can replay it byte-for-byte.

\subsection{Stage 2 -- Cross-validated evaluation}
\label{sec:method-cv}
Each design row is evaluated under stratified $k$-fold CV
(default $k{=}5$); per-fold metrics are persisted in a long-format
CSV. Quality metrics include accuracy, precision, recall, specificity,
F1, ROC-AUC, PR-AUC, balanced accuracy, MCC, Brier score, and
expected calibration error (ECE). Wall-clock cost is the mean
fit-plus-predict time per fold. The single source of truth for
hyperparameter naming and types is \texttt{doe\_xgb.config}.

\subsection{Stage 3 -- Factor model with Varimax rotation}
\label{sec:method-fa}
The article-track factor model
(\texttt{doe\_xgb.factor\_model}) supports four modes:
\textsc{auto} (Kaiser eigenvalue $>1$ \emph{or} cumulative variance
threshold), \textsc{fixed} (explicit number of factors -- the
article default for the headline run is three),
\textsc{manual} (user-declared constructs such as
Quality / Cost / Robustness), and \textsc{none}
(no FA; raw metrics flow through unchanged). All modes return
loadings, factor scores, eigenvalues, explained variance,
communalities, the rotation matrix, the construct map, and a KMO
estimate where defined.

\subsection{Stage 4 -- Objective specification}
\label{sec:method-objectives}
Every objective entering the optimization is declared through an
\texttt{ObjectiveSpec} (\texttt{doe\_xgb.objectives}) with required
fields \texttt{direction}
$\in \{\textsc{minimize},\textsc{maximize},\textsc{target}\}$,
\texttt{transform}
$\in \{\textsc{raw},\textsc{zscore},\textsc{log1p},\textsc{factor\_score},\textsc{fmse},\textsc{custom}\}$,
\texttt{role}
$\in \{\textsc{primary\_nbi},\textsc{post\_optimization},\textsc{reporting},\textsc{constraint}\}$,
and optional \texttt{target}, \texttt{weight}, \texttt{group},
\texttt{bounds}, and \texttt{normalization}.
A canonicalizer maps every spec to a minimization callable. For
\textsc{target} / FMSE objectives we use Eq.~21 of
\citet{pereira2025eaai}:
\begin{equation}
  f^{\text{FMSE}}_i(\bm x) \;=\; \bigl(\hat F_i(\bm x) - T_i\bigr)^{2}
                                  \;+\;\sigma^{2}_{\hat F_i},
  \label{eq:fmse}
\end{equation}
where $T_i$ is the target value (e.g., the per-objective surrogate
optimum) and $\sigma^{2}_{\hat F_i}$ is the factor-variance term
implied by the loadings $\bm L$ of stage 3. The article-track
default for $\VRF$ objectives is \textsc{target} + FMSE; raw ML
metrics (e.g., \emph{Time\_MeanFold}, \emph{Accuracy\_Mean}) use
\textsc{raw} with explicit \textsc{minimize} / \textsc{maximize}.

\subsection{Stage 5 -- Surrogate fitting}
\label{sec:method-surrogate}
For process variables (CCDFC, Box--Behnken, factorial), we fit a
full quadratic RSM with strong-hierarchy backward elimination at
$\alpha{=}0.05$ (\texttt{ProcessQuadraticRSM}). For mixture variables
on the simplex (used by the post-optimization stage of
Section~\ref{sec:method-mbpa}), we fit a Scheff\'{e} canonical
mixture polynomial of order linear / quadratic / special cubic / cubic
(\texttt{MixtureScheffeModel}). Backward elimination is intentionally
disabled for Scheff\'{e} models because the canonical interpretation
depends on the full term set. A
\texttt{validate\_for\_model} routine warns or rejects incompatible
design/family combinations (mixture model on CCDFC; process quadratic
on simplex; quadratic on Latin hypercube).

\subsection{Stage 6 -- True $N$-objective NBI}
\label{sec:method-nbi}
Let $\hat{\bm F}(\bm x) = (\hat F_1(\bm x),\dots,\hat F_q(\bm x))$ be
the canonicalized minimization surrogates. We compute per-objective
anchors $\bm x^{*}_j = \arg\min_{\bm x\in\Omega} \hat F_j(\bm x)$ and
evaluate every objective at every anchor to obtain the payoff matrix
$\bm{\Phi}^{\text{full}}_{ij} = \hat F_i(\bm x^{*}_j)$. The utopia
point is its diagonal,
$\bm F^{U}_i = \hat F_i(\bm x^{*}_i)$, and the CHIM matrix is
$\bm{\Phi} = \bm{\Phi}^{\text{full}} - \bm F^{U}\bm{1}^{\top}$. The
quasi-normal direction defaults to the Das--Dennis convention,
\begin{equation}
  \hat{\bm n} \;=\; \dfrac{-\bm{\Phi}\,\bm{1}}{\|\bm{\Phi}\,\bm{1}\|},
  \label{eq:nhat}
\end{equation}
with a Gram--Schmidt orthogonal-complement variant available as a
diagnostic. For a weight vector $\bm\beta$ on the simplex the
sub-problem is
\begin{equation}
  \begin{aligned}
    \max_{\bm x \in \Omega,\, t \ge 0} \quad & t \\
    \text{s.t.}\quad & \hat{\bm F}(\bm x) - \bm F^{U}
       - \bm{\Phi}\bm\beta - t\,\hat{\bm n} = \bm 0,
  \end{aligned}
  \label{eq:nbi-subproblem}
\end{equation}
solved as $q$ scalar equality constraints and $\bm x$ box bounds via
SLSQP with multistart. Weight vectors come from a Scheff\'{e}
simplex-lattice $\{q,m\}$ (default $m{=}50$ for $q{=}2$, $m{=}10$ for
$q{=}3$, giving 51 and 66 sub-problems respectively). The residual
$\|\hat{\bm F}(\bm x^{*}(\bm\beta)) - \bm F^{U} - \bm{\Phi}\bm\beta -
t\,\hat{\bm n}\|$ is reported for every sub-problem so reviewers can
audit constraint satisfaction.

\paragraph{Methodology pin.} The $N$-objective generality is
verified by automated tests for $q\in\{2,3,4,5\}$ on synthetic
quadratics with closed-form anchors; residual norms remain below
$10^{-3}$ across the entire test matrix.

\subsection{Stage 7 -- Conditional MBPA post-optimization}
\label{sec:method-mbpa}
A flat or compressed primary frontier reduces the practical value of
$\NBI$. We compute six diagnostics on the front -- normalized
objective range, average pairwise distance, count of unique
non-dominated points, weight concentration, curvature score, spread
$\Delta$ -- and trigger MBPA when any one fails its threshold
(\texttt{doe\_xgb.diagnostics}). When triggered
(\texttt{doe\_xgb.post\_optimization}) we fit Scheff\'{e} canonical
mixture surrogates over the simplex of weights for the
generalized distance to utopia
\citep[Eq.~29]{pereira2025eaai} and the Shannon entropy
$S(\bm w) = -\sum_i w_i \ln w_i$, then minimize $\mathrm{GD}(\bm w) -
S(\bm w)$ inside an elliptical interior constraint that excludes
vertex weights. Default radii are $0.30$, matching the EAAI
recommendation. When MBPA is skipped we still write
\texttt{frontier\_quality.json} so the decision is auditable.

\subsection{Final selection}
\label{sec:method-selection}
A pluggable selection rule (\texttt{doe\_xgb.selection}) chooses the
final candidate. The article-track default is
\textsc{distance\_to\_utopia}; \textsc{knee},
\textsc{utility}, \textsc{lexicographic}, and \textsc{topsis} are
available. Benchmarks retain \textsc{max\_accuracy} so the
dissertation tables remain comparable. The full Pareto front is
always persisted regardless of the selection rule.

\subsection{Legacy weighted-sum ablation}
\label{sec:method-legacy}
The dissertation-era weighted-sum scalarization
(\texttt{doe\_xgb.scalarization.run\_nbi\_weighted\_sum}) is preserved
verbatim and exposed as an ablation baseline. It is never referred to
as NBI in the article text or in the article-track code; the
discrepancy is documented in
\texttt{docs/METHODOLOGY\_DECISIONS.md} D1.

\section{Experimental design}\label{sec:experimental-design}

\subsection{Primary benchmark: OpenML-CC18}
\label{sec:exp-datasets}
The primary doctoral benchmark is the OpenML-CC18 curated
classification suite of \citet{bischl2021openmlbenchmark}
(\texttt{suite\_id}~$=$~99), containing 72 standardized classification
tasks served by the OpenML platform \citep{vanschoren2014openml}. The
suite was constructed by filtering OpenML for tasks with verified
licences, well-defined train/test splits, dataset sizes between
500 and 100\,000 rows (with a small number of larger exceptions),
$\geq\!2$ classes, and reasonable class balance, and is widely used
as a fair common ground for HPO and AutoML evaluation. We use it as
\emph{a task-based benchmark, not a dataset-based one}: a few CC18
tasks share the same underlying OpenML dataset under different target
columns or train/test splits, so identity is by
\texttt{openml\_task\_id} and never by dataset name alone. The
materialized task / dataset registry, suite metadata, and a
binary / multiclass / categorical / imbalance / size-bucket coverage
report are committed under
\texttt{benchmarks/doctoral/openml\_cc18/} and are produced
deterministically by \texttt{scripts/import\_openml\_cc18.py}.

\paragraph{Coverage of the panel.} The 72 tasks split as 35 binary
and 37 multiclass; 24 tasks have at least one categorical feature;
19 tasks have a class-imbalance ratio $\geq\!5$. By size, 18 tasks
have $\leq\!1\,000$ rows, 43 have $1\,001$--$30\,000$, and 11 have
$>\!30\,000$. The full per-task table (rows, features, $n$ classes,
class distribution, license, OpenML URL) lives in
\texttt{benchmarks/doctoral/openml\_cc18/tasks.csv}; we do not
reproduce it here because the suite is canonical and externally
specified.

\paragraph{Internal smoke panel (not part of the benchmark).} The
12-dataset panel that drove earlier integration testing
(\texttt{magic}, \texttt{breast\_cancer}, \texttt{pima\_diabetes},
\texttt{spambase}, \texttt{adult}, \texttt{bank\_marketing},
\texttt{credit\_card\_default}, \texttt{german\_credit},
\texttt{wine\_quality}, \texttt{dry\_bean}, \texttt{mushroom},
\texttt{phishing}) is preserved as a development / smoke / runtime
profiling fixture under
\texttt{benchmarks/doctoral/internal\_smoke\_panel/} and is
\emph{not} reported as part of the doctoral benchmark.

\subsection{Algorithms}
\label{sec:exp-algorithms}
We compare three GBDT families: XGBoost \citep{chen2016xgboost},
LightGBM \citep{ke2017lightgbm}, and CatBoost
\citep{prokhorenkova2018catboost}. Categorical features are passed
natively to CatBoost; for XGBoost and LightGBM the headline run uses
label-encoded ints for cross-method parity with HPO baselines that
do not understand categorical columns natively (this default is
toggleable per algorithm and is logged in the per-replica manifest).
The proposed method and every baseline tune the same per-algorithm
hyperparameter list under the same bounds, listed in
Table~\ref{tab:hpo-spaces}.
\todoblock{Confirm the LightGBM / CatBoost search spaces. Initial
plan: keep the seven-factor XGBoost space for parity with the
dissertation, and use analogues for LightGBM (\texttt{num\_leaves},
\texttt{min\_data\_in\_leaf}) and CatBoost (\texttt{depth},
\texttt{l2\_leaf\_reg}).}

\begin{table}[t]
  \centering
  \caption{Per-algorithm hyperparameter search spaces (tentative).}
  \label{tab:hpo-spaces}
  \begin{tabular}{lll}
    \toprule
    Algorithm & Hyperparameter & Range \\
    \midrule
    \multirow{7}{*}{XGBoost}
      & \texttt{subsample}         & $[0.05, 1.00]$ \\
      & \texttt{colsample\_bytree} & $[0.05, 1.00]$ \\
      & \texttt{colsample\_bylevel}& $[0.05, 1.00]$ \\
      & \texttt{learning\_rate}    & $[0.01, 0.30]$ \\
      & \texttt{max\_depth}        & $\{3,\dots,18\}$ \\
      & \texttt{gamma}             & $[0.05, 5.00]$ \\
      & \texttt{n\_estimators}     & $\{50,\dots,700\}$ \\
    \midrule
    \multirow{7}{*}{LightGBM} & \multicolumn{2}{l}{\todoblock{fill in.}}\\
    \midrule
    \multirow{7}{*}{CatBoost} & \multicolumn{2}{l}{\todoblock{fill in.}}\\
    \bottomrule
  \end{tabular}
\end{table}

\subsection{Comparative protocol and method matrix}
\label{sec:exp-protocol}
The methods entering the doctoral benchmark are frozen by the
machine-readable matrix at
\texttt{benchmarks/doctoral/openml\_cc18/method\_matrix.csv}; the
authoritative narrative specification is
\texttt{docs/COMPARATIVE\_PROTOCOL.md}. The matrix is the single
source of truth for shard generation: the script that materializes
the SQLite job matrix reads its method list from that CSV and never
hardcodes method names. Table~\ref{tab:methods} summarizes the
matrix. Methods carry a \emph{primary / subset / ablation /
literature\_only} status and an \emph{objective mode} flag (single-
versus multi-objective).

\begin{table}[t]
  \centering
  \caption{Methods participating in the OpenML-CC18 doctoral
  benchmark. ``Status'' is one of \emph{primary} (runs on all 72
  tasks at every stage), \emph{subset} (runs only on the CC18 subset
  defined in \texttt{benchmarks/doctoral/openml\_cc18/comparative\_protocol.md}),
  \emph{ablation} (joins from stage 2 onward), or \emph{literature}
  (cited as context, not benchmarked). Multi-fidelity rows share a
  budget-equivalence rule based on total boosting iterations rather
  than configuration count. The full machine-readable specification
  is \texttt{benchmarks/doctoral/openml\_cc18/method\_matrix.csv}.}
  \label{tab:methods}
  \small
  \begin{tabular}{lllll}
    \toprule
    Method & Family & Status & Objective & Reference \\
    \midrule
    \texttt{default\_gbdt}                & control          & primary    & single & --- \\
    \texttt{random\_search}               & classical        & primary    & single & \citep{bergstra2012random} \\
    \texttt{tpe\_optuna}                  & classical        & primary    & single & \citep{bergstra2011tpe,akiba2019optuna} \\
    \texttt{smac3}                        & SMAC             & primary    & single & \citep{hutter2011smac,lindauer2022smac3} \\
    \texttt{hyperband\_or\_asha}          & multi-fidelity   & primary    & single & \citep{li2017hyperband,li2020asha} \\
    \texttt{bohb}                         & multi-fidelity   & primary    & single & \citep{falkner2018bohb} \\
    \texttt{dehb}                         & multi-fidelity   & primary    & single & \citep{awad2021dehb} \\
    \texttt{nsga2}                        & evolutionary MOO & primary    & multi  & \citep{deb2002nsga2,blank2020pymoo} \\
    \texttt{motpe}                        & MOO              & primary    & multi  & \citep{ozaki2020motpe} \\
    \texttt{parego}                       & MOO              & subset     & multi  & \citep{knowles2006parego} \\
    \texttt{doe\_rsm\_vrf\_true\_nbi}     & proposed         & primary    & multi  & this work \\
    \texttt{\ldots\_no\_mbpa}             & ablation         & ablation   & multi  & this work \\
    \texttt{legacy\_weighted\_sum}        & ablation         & ablation   & multi  & \citep{ribeiro2026dissertation} \\
    \texttt{flaml\_optional}              & AutoML context   & literature & single & \citep{wang2021flaml} \\
    \bottomrule
  \end{tabular}
\end{table}

\paragraph{Replicas and CV.} Each (task, algorithm, method) cell is
repeated for $R{=}30$ replicas, staged $1\!\to\!5\!\to\!10\!\to\!30$
through the four stages \texttt{stage0\_replica\_001},
\texttt{stage1\_topup\_to\_005}, \texttt{stage2\_topup\_to\_010},
\texttt{stage3\_topup\_to\_030}. Each replica fits the OpenML
task-defined CV folds when available and otherwise stratified
5-fold. Replica seeds are derived from a single \texttt{seed\_base}
(default 42) so the protocol can be replayed byte-for-byte.

\paragraph{Budget equivalence.} Single-objective + multi-objective
non-fidelity methods all run at the same headline budget
$B = \mathrm{DOE\_RUNS} + \mathrm{NBI\_EVAL\_K} = 88 + 50 = 138$
configurations per replica per (task, algorithm), each evaluated
under the same CV split. Multi-fidelity methods are budgeted by
total boosting iterations rather than configuration count
($\text{total\_n\_estimators\_budget} = B \cdot \texttt{max\_iter}$).
Evolutionary methods set population size and generations so
$\text{pop\_size} \cdot \text{gens} = B$. The full specification is in
Section~``Budget equivalence rule'' of
\texttt{docs/COMPARATIVE\_PROTOCOL.md}.

\paragraph{Headline target.} For the methods that run on all 72 CC18
tasks (every row of the method matrix with \texttt{full\_cc18}~$=$
\texttt{true}), the campaign is

\[
  72\ \text{tasks} \times 3\ \text{algorithms} \times
  n_{\text{methods,full}} \times 30\ \text{replicas},
\]
plus the per-replica jobs contributed by the \emph{subset} methods on
their defined CC18 subset. The exact $n_{\text{methods,full}}$ is
counted from \texttt{method\_matrix.csv} at shard-generation time and
is not duplicated as a constant in the manuscript.

\subsection{Metrics and statistical analysis}
\label{sec:exp-metrics}
\paragraph{Quality metrics.} ROC-AUC, PR-AUC, F1 (binary) /
F1-macro (multiclass), balanced accuracy, MCC. Multiclass tasks use
the one-vs-rest variants of the AUC metrics.

\paragraph{Calibration metrics.} Brier score (binary or multiclass)
and Expected Calibration Error (ECE)
\citep{naeini2015calibration}.

\paragraph{Cost metrics.} Mean fit-plus-predict time per fold;
total optimization time; total wall-clock time per replica.

\paragraph{Multi-objective frontier diagnostics.} For
multi-objective methods we additionally report: number of unique
non-dominated points; spread $\Delta$; weight-concentration index;
NBI sub-problem residual statistics (where applicable); MBPA-trigger
state. These are emitted regardless of whether MBPA fired so that
the conditional decision is always auditable.

\paragraph{Statistical analysis.} Paired comparisons within each task
use a non-parametric Wilcoxon signed-rank test (or paired $t$ after
log-transform when the residuals are reasonably Gaussian). Across
tasks, we use the Friedman test followed by Nemenyi post-hoc with
critical-difference plots, following
\citet{demsar2006statistical}. All $p$-values use Holm--Bonferroni
correction within each family of comparisons. Effect sizes are
reported as the median of paired differences plus a non-parametric
bootstrap 95\% confidence interval; rank statistics are computed
both per task and aggregated across the 72 tasks.

\paragraph{Robustness across replicas.} For each (task, algorithm,
method) cell we report the across-replica IQR of every quality
metric, and we additionally report the rate of replica-level
failures (job exits with status \texttt{failed}) so that the
robustness of each method is visible at the task level.

\subsection{Cost discipline and reproducibility}
\label{sec:exp-cost-discipline}
\paragraph{Cost discipline.} Before any shard generation,
\texttt{doe-xgb estimate-cost --preset
openml\_cc18\_72\_tasks\_3\_algorithms\_30\_replicas} produces an
authoritative wall-clock estimate. The campaign is gated by a
sign-off on this estimate.

\paragraph{Determinism.} Headline tables use
\texttt{tree\_method=exact} and \texttt{n\_jobs=1} (article-track
default); the determinism cost is documented in
\texttt{docs/METHODOLOGY\_DECISIONS.md} D13.

\paragraph{Sharded execution.} The job matrix is sharded across
SQLite files under \texttt{jobs/doctoral/openml\_cc18/shards/} so
that the dedicated MacBook Pro can run the campaign with low lock
contention; the schema (Section~\ref{sec:method-infra}) tracks
claim / success / failure transitions and per-job runtime so the
campaign can resume after interruption without re-running successful
jobs.

\paragraph{Reproducibility.} Every replica writes a
system-fingerprinted manifest, the full Pareto front (for
multi-objective methods), the NBI sub-problem residuals
(for the proposed method), the MBPA trigger diagnostics, and a
long-format per-fold metric CSV. Aggregation scripts regenerate
Tables~\ref{tab:performance}--\ref{tab:cost} of
Section~\ref{sec:results} from those artifacts.

\section{Results}\label{sec:results}

\todoblock{This section is intentionally a scaffold. Insert results
\emph{verbatim} from the aggregation CSVs once
\texttt{configs/article\_3vrf\_xgb\_magic.yaml} (and its multi-dataset
variant to be created) finishes running on the v1 panel of
Section~\ref{sec:experimental-design}. Do not invent numbers.}

\subsection{Dataset-level performance}
\label{sec:res-performance}
\begin{table}[t]
  \centering
  \caption{Predictive performance (mean $\pm$ standard deviation over
  $R{=}10$ replicas) per (dataset, algorithm, method).
  \textbf{TODO: insert after experiment runs.}}
  \label{tab:performance}
  \begin{tabular}{lllrrrrr}
    \toprule
    Dataset & Algorithm & Method & ROC-AUC & PR-AUC & F1 & Bal.\ Acc. & MCC \\
    \midrule
    \multicolumn{8}{c}{\textit{TODO: rows from
       experiments/<dataset>/<design>/aggregated\_table.csv}} \\
    \bottomrule
  \end{tabular}
\end{table}

\subsection{Computational cost}
\label{sec:res-cost}
\begin{table}[t]
  \centering
  \caption{Computational cost: mean time per fold, total optimization
  time, and total wall-clock time. \textbf{TODO}.}
  \label{tab:cost}
  \begin{tabular}{lllrrr}
    \toprule
    Dataset & Algorithm & Method & Time/fold (s) & Opt.\ time (s) & Total time (s) \\
    \midrule
    \multicolumn{6}{c}{\textit{TODO}} \\
    \bottomrule
  \end{tabular}
\end{table}

\subsection{Cross-dataset comparison}
\label{sec:res-friedman}
\todoblock{Insert Friedman test and Nemenyi critical-difference plot
after running the panel. Discuss any rank inversions between the
proposed method and TPE/BO. Cite \citet{demsar2006statistical}.}

\subsection{NBI vs.\ legacy weighted-sum ablation}
\label{sec:res-ablation}
\todoblock{Side-by-side comparison: legacy
\texttt{run\_nbi\_weighted\_sum} vs.\ true $N$-objective NBI. The
expected outcome is that NBI produces more uniformly spaced points on
the Pareto front, but the practical performance difference depends on
front curvature; a synthetic confirmation already passes in
\texttt{tests/methodology/test\_nbi\_core.py}.}

\subsection{$N$-objective NBI residual diagnostics}
\label{sec:res-nbi-residuals}
\begin{table}[t]
  \centering
  \caption{NBI sub-problem residual norms across the $51$ $\beta$ /
  $66$ $\beta$ grid for the $q{=}2$ / $q{=}3$ runs, per dataset.
  \textbf{TODO}.}
  \label{tab:nbi-residuals}
  \begin{tabular}{lrrr}
    \toprule
    Dataset & $\max \|\text{residual}\|$ & median & $p_{95}$ \\
    \midrule
    \multicolumn{4}{c}{\textit{TODO}} \\
    \bottomrule
  \end{tabular}
\end{table}

\subsection{Factor model and VRF interpretation}
\label{sec:res-vrf}
\todoblock{Insert per-dataset rotated loading tables (one row per raw
metric, one column per VRF), with eigenvalues and cumulative
variance. Highlight the construct each VRF maps to.}

\subsection{MBPA post-optimization trigger summary}
\label{sec:res-mbpa}
\begin{table}[t]
  \centering
  \caption{Per-dataset MBPA trigger summary. ``Triggered'' = at least
  one of the six diagnostics fired; ``Refined'' = a refined
  $\bm w^{*}$ was recovered; ``Improved'' = the refined candidate
  improved the selection metric over the primary $\bm w^{*}$.
  \textbf{TODO}.}
  \label{tab:mbpa}
  \begin{tabular}{lrrrr}
    \toprule
    Dataset & Replicas & Triggered & Refined & Improved \\
    \midrule
    \multicolumn{5}{c}{\textit{TODO}} \\
    \bottomrule
  \end{tabular}
\end{table}

\subsection{Reproducibility artifact summary}
\label{sec:res-artifacts}
\todoblock{Insert artifact count, total storage, hash of the canonical
design, and the GitHub release tag (\texttt{v0.3.0-articlev1}) once
the campaign is archived.}

\section{Discussion}\label{sec:discussion}

\todoblock{Defer the discussion until the v1 results are aggregated.
Below is a checklist of points the discussion must engage with so that
the article remains honest about scope and limitations.}

\paragraph{Why true NBI matters in practice.} The dissertation
implementation \citep{ribeiro2026dissertation} produced compelling
cost-quality numbers using a weighted-sum scalarization that the
post-defense audit reclassified as \emph{not NBI}.
\todoblock{Compare cost-quality outcomes between the legacy
weighted-sum and the article-track $N$-objective NBI on the same v1
panel. If parity holds across most datasets but NBI dominates on
non-convex fronts, that is the headline contribution.}

\paragraph{When MBPA fires.}
The conditional trigger fires only when the primary frontier is
flat / compressed. \todoblock{Report the trigger frequency per
dataset and per algorithm; confirm the 0.30 elliptical radius default
or recommend an updated value.}

\paragraph{When the framework over-engineers the problem.}
Small datasets (Breast Cancer, German Credit) and high-separability
ones (Mushroom) may collapse to nearly flat fronts at near-perfect
quality; in that regime, classic random search with a tighter budget
may be the right answer. \todoblock{Quantify and discuss.}

\paragraph{Algorithm-specific findings.}
\todoblock{Report any cases where the proposed method is markedly
better / worse for CatBoost than for XGBoost / LightGBM; relate to
CatBoost's ordered boosting and categorical handling.}

\paragraph{Limitations.}
The v1 panel is restricted to twelve public tabular datasets and ten
replicas; statistical claims are necessarily within that scope. The
calibration metrics (Brier, ECE) are reported but not optimized; the
proposed method is tuning for quality and cost, and calibration is a
secondary lens.
\todoblock{Add limitations specific to the observed results --
e.g., MBPA trigger rate, anchor convergence on certain datasets,
sensitivity to the choice of factor count.}

\paragraph{Threats to validity.}
\todoblock{Discuss: (i) the same seed protocol yields tiny but
non-zero non-determinism on \texttt{tree\_method=hist}; the article
defaults to \texttt{exact} for the headline tables; (ii) the cost
estimator over-estimates by the configured \texttt{overhead\_factor}
of 1.10; (iii) the 51 / 66 NBI sub-problem grid may under-cover
non-convex fronts in 4D+; the slow-marked test for $q{=}5$ runs only
the centroid.}

\paragraph{Implications for the doctoral 82-dataset benchmark.}
\todoblock{Use the cost estimator to project from the v1 results to
the doctoral scope. State explicitly that the present article is
\emph{not} the doctoral benchmark.}

\section{Conclusion}\label{sec:conclusion}

We presented a post-dissertation methodological refinement and
generalization of the DoE + RSM + Varimax + NBI framework for
multi-objective hyperparameter tuning of GBDT models. The article-track
implementation departs from the dissertation pipeline of
\citep{ribeiro2026dissertation} in four substantive ways:
(i) the optimizer is a true $N$-objective NBI rather than a
weighted-sum scalarization; (ii) every objective carries an explicit
direction, with FMSE wrapping for $\VRF$ targets; (iii) the surrogate
fitting is design-aware, separating process-quadratic RSM from
Scheff\'{e} mixture polynomials; and (iv) a conditional MBPA
post-optimization stage refines weights when the primary frontier is
weakly informative. The legacy weighted-sum solver is retained as an
ablation baseline and is no longer named NBI in either the code base
or the article text.

\todoblock{Quantitative claims (CPU-hour reduction, predictive parity
intervals, MBPA trigger rate, NBI residual percentiles) go here once
the v1 campaign is aggregated. Avoid hyperbole; report medians and
inter-quartile ranges, not single-replica winners.}

\paragraph{Future work.} Three directions follow naturally:
the doctoral 82-dataset benchmark; an extension to non-tree
estimators with the same \texttt{ObjectiveSpec} layer
(linear models, deep tabular networks); and a tighter coupling
between the cost estimator and the orchestrator so that a campaign
can self-tune its replica count to a wall-clock budget.


% --- Acknowledgements --------------------------------------------------------
\section*{Acknowledgements}
\todoblock{Add funding agencies, advisors, computational resources, and
co-authors. Mention UNIFEI, FAPEMIG (BPD-01045-22 collaboration if
applicable), and the GitHub artifact / CI infrastructure hosting the
reproducibility bundle.}

% --- Code/data availability ---------------------------------------------------
\section*{Code and data availability}
The article-track implementation, configuration files, and
reproducibility manifests are openly available under the MIT license at
\url{https://github.com/caioribeiro99/doe_nbi_hpo_project} on the
\texttt{repo-publication-readiness} branch. The dissertation baseline is
preserved on \texttt{main}, tagged \texttt{v0.1.0-dissertation}.
Every replica writes a system-fingerprinted manifest and per-fold
metrics; aggregation scripts regenerate the tables in
Section~\ref{sec:results}.

% --- Bibliography -------------------------------------------------------------
\bibliographystyle{plainnat}
\bibliography{references}

\end{document}
